\section{Open problems}\label{sec:open}

\begin{enumerate}[(1),leftmargin=*]
\item \emph{The order of the threshold.} Theorem~\ref{thm:threshold} gives
$q_h\le m_h\le h^{2+o(1)}$. Is $m_h=O(h^2)$? The known values are $m_4=6$, $m_5=8$, and
$m_6=15$, and $m_7\le21$ (Remark~\ref{rem:seven}). What is $m_7$, and what is $f(7,m)$ for
$9\le m\le20$?
\item \emph{A lower bound.} The bound $m_h\ge q_h$ only counts the doubles. Is there a
constant $c>1$ with $m_h\ge c\,q_h$ for all large $h$? Does $m_h/q_h$ tend to infinity?
\item \emph{Monotonicity.} Is $f(h,m+1)\ge f(h,m)+1$ for all $h\ge2$ and $m\ge3$? It would
imply that the widths with $f(h,m)=m+q_h$ are exactly the widths $m\ge m_h$. We could not even prove that $f(h,m)$ is
nondecreasing in $m$.
\item \emph{Square grids.} What is the order of growth of $f(n,n)$, between $n^{6/5-o(1)}$
and $n^{4/3}$? Theorem~\ref{thm:small} gives $f(n,n)\le2n$ for $n\le7$. What is the first
$n$ with $f(n,n)>2n$?
\end{enumerate}
